\documentclass[conference]{IEEEtran}
\usepackage[utf8]{inputenc}
\usepackage{graphicx,amsmath,amssymb,booktabs,hyperref,url,xcolor,listings,enumitem}
\hypersetup{colorlinks=true,linkcolor=blue,citecolor=blue,urlcolor=blue,breaklinks=true}
\lstset{basicstyle=\footnotesize\ttfamily,breaklines=true,frame=single,language=Python}
% Fast fix for overflow without heavy microtype
\sloppy
\setlength{\emergencystretch}{2.5em}
\title{Failures: Deterministic Failure-Mode Guardrails Make Coding Agents Build Resilient Systems}
\author{\IEEEauthorblockN{Mayowa Kalejaiye}
\IEEEauthorblockA{Failures MCP Project\\ \texttt{github.com/mayowa-kalejaiye/Failures}}}
\begin{document}
\maketitle

\begin{abstract}
Coding agents write happy-path code because their training data is happy-path tutorials. We present \textbf{Failures}, a deterministic Model Context Protocol (MCP) server that forces agents to reason about failure modes \emph{before code ships}---via 11 engineering invariants (atomicity, idempotency, timeout/ambiguous outcome, concurrency, ordering, consistency, availability, resource exhaustion, recovery, observability, retry safety) checked by static rules with line evidence and calibrated confidence, not LLM prompts. Zero tokens, zero hallucination. Across 6 realistic scenarios (payment via Paystack, queue, authentication, file upload, inventory, webhook) with blind evaluation (same prompt without vs.\ with Failures, 3 agents: Claude, Cursor/Codex, scored by a deterministic harness on 10 resilience criteria), Failures-enabled agents eliminate CRITICAL findings (payment 3$\to$0, webhook 3$\to$0) and surface the correct state machine (\texttt{pending$\to$Paystack$\to$ambiguous$\to$reconciliation$\to$verified$\to$enrollment}) without being asked. We also contribute 6 adversarial \emph{looks-safe-but-isn't} cases and a proportional-architecture metric to penalize over-engineering. All artifacts, prompts, and manifests (\texttt{prompt\_hash}, commit) are frozen for reproducibility.
\end{abstract}
\begin{IEEEkeywords}AI-assisted programming, resilience, distributed systems, MCP, static analysis, empirical software engineering\end{IEEEkeywords}

\section{Introduction}
Modern coding agents can scaffold a payment system in minutes. But ask them to ``Build a course payment system using Paystack. After payment, automatically enroll the student''---and they will happily write \texttt{charge(); enroll();} with no idempotency key, no transaction boundary, no handling for the most dangerous case: \emph{ambiguous outcome} (provider succeeded, but the response was lost). In production this is a double charge or a lost enrollment.

The problem is structural: agents are trained on tutorials, not incident reports. Resilience knowledge---idempotency, outbox, reconciliation, optimistic locking---lives in books like \emph{Release It!}~\cite{releaseit} and \emph{DDIA}~\cite{ddia}, not in the prompt.

\paragraph{Insight.} Resilience can be checked \emph{deterministically} before code runs. We encode 11 failure dimensions as testable invariants with static rules that return line evidence and honest confidence, never ``BUG PROVEN'' unless syntactically provable. An MCP server exposes these as 13 tools; the agent is required to call \texttt{review\_plan} and \texttt{review\_code} before finalizing.

\paragraph{Contributions.}
(1) A failure model and specification (v0.2.0) covering hard/ambiguous/partial/duplicate failures and 11 dimensions with a Principle$\rightarrow$FailureMode$\rightarrow$Pattern$\rightarrow$Test graph.
(2) \textbf{Failures MCP} 0.3.1---deterministic, zero-token checks with evidence/confidence, live on PyPI (\texttt{pipx install failures-mcp}).
(3) A blind evaluation harness over 6 scenarios + 6 adversarial ``looks-safe'' cases, with a 10-criterion resilience score and a proportional-architecture guard against over-engineering.
(4) Full reproducibility: frozen prompts (\texttt{prompt\_hash}), manifests, commits, and agent outputs for Claude, Cursor (Grok 4.6), and Codex.

\paragraph{Result in one line.} Same prompt, same agent, with Failures connected: payment queue, inventory, and webhook systems go from CRITICAL-risk to zero CRITICAL, and the correct \texttt{pending+idempotency+webhook dedup+reconciliation} architecture emerges unprompted.

\section{Related Work}
\paragraph{AI code generation evaluation.} Copilot and agent benchmarks focus on functional correctness (HumanEval, SWE-bench)~\cite{copilot-eval}. We shift the metric to \emph{final system resilience} under failures, not pass@k.

\paragraph{Resilience patterns.} Nygard~\cite{releaseit}, Kleppmann~\cite{ddia}, and SRE~\cite{sre} catalog patterns (circuit breaker, outbox, idempotency-key). Failures makes them \emph{machine-checkable} via invariants and ties each finding to a required pattern and test.

\paragraph{Static guards vs.\ LLM judges.} LLM-as-judge approaches hallucinate and cost tokens. Our checks are regex/AST heuristics labeled \texttt{STATIC FAILURE CHECK---Potential violation. Confidence: X}. High-confidence rules ($>$0.95) are syntactic proofs (e.g., two \texttt{execute} without \texttt{BEGIN}); weaker signals are surfaced as Medium with evidence, never as proven bugs.

\section{Failure Model and Specification}
We follow \texttt{FAILURES\_SPEC.md} v0.2.0, the single source of truth.

\paragraph{Failure taxonomy.} Hard failure (500/exception), \emph{ambiguous outcome} (caller does not know if side effect happened---timeout, crash before ack, partition), partial success, duplicate execution (retry/redelivery).

\paragraph{Principle.} A named, testable invariant: \{id, question, invariant, applies\_to, failure\_modes, patterns, tests\}. Example---idempotency: ``What happens if this operation executes more than once?'' Invariant: one logical operation yields at most one effect.

\paragraph{11 dimensions.} atomicity, idempotency, ordering, concurrency, availability, timeout, consistency, resource\_exhaustion, recovery, observability, retry\_safety. Each defines a builder question (e.g., timeout: ``What happens when you don't know whether it succeeded?'').

\paragraph{Finding.} Unit of output: \{id, dimension, severity (CRITICAL/HIGH/MEDIUM/LOW), confidence 0--1, mode, title, why, evidence\{excerpt, lines\}, risk, required/optional mitigations, tests\}. Severity is impact-anchored; confidence is honest.

\paragraph{Invariants.} Predicates that must hold under \emph{any} failure, e.g., \texttt{invariant.payment.not\_double\_charged}, \texttt{invariant.enrollment.requires\_successful\_payment}. \texttt{check\_invariant} returns violatable? + counter-scenario.

\section{Approach: Deterministic MCP Guardrails}
\paragraph{Architecture.} Knowledge graph \texttt{Principle$\rightarrow$FailureModes$\rightarrow$Pattern$\rightarrow$Test} materialized as \texttt{mcp\_server/knowledge/\{principles,rules,dimensions\}.json}. All 13 tools are deterministic-first; LLM is optional post-processing only.

\begin{table}[h]\centering\small
\begin{tabular}{ll}
\toprule Tool & Purpose \\
\midrule
\texttt{get\_principles} & list principles/dimensions \\
\texttt{review\_architecture} & findings from system/components/flows \\
\texttt{analyze\_component} & component-specific findings \\
\texttt{review\_plan} & plan-step findings + reordering \\
\texttt{check\_invariant} & violatable invariants + counter-scenario \\
\texttt{generate\_failure\_cases} & concrete cases \\
\texttt{review\_code} & findings with evidence/confidence \\
\texttt{generate\_failure\_tests} & failure tests \\
\texttt{check\_idempotency/retry/transaction} & decision + required controls \\
\texttt{list\_failures} & rules + docs \\
\bottomrule
\end{tabular}
\caption{Failures MCP tools (abridged). Outputs are severity-sorted, with \texttt{--- JSON ---} machine block.}
\end{table}

\paragraph{Example.} \texttt{review\_code} on \texttt{await stripe.charge(...)} without \texttt{idempotency\_key} yields CRITICAL \texttt{external\_call\_without\_idempotency} (confidence 0.91, evidence \texttt{42-43: no idempotency\_key}), risk ``Retry after ambiguous timeout may duplicate charge,'' required ``persist key BEFORE call, UNIQUE constraint, reconciliation worker,'' and tests \texttt{retry\_after\_provider\_success, duplicate\_request}.

\paragraph{Design choice.} Checks emit \texttt{STATIC FAILURE CHECK} with confidence, never ``bug proven'' unless syntactic. This preserves trust and auditability, and keeps cost at zero tokens.

\section{Evaluation Design}
We answer: \emph{Does connecting Failures to a coding agent measurably improve the failure resilience of the software it produces?} Metric is \emph{final system}, not finding count.

\paragraph{Scenarios (6).} \texttt{scenarios/*.md} deliberately omit resilience hints:
\texttt{payment} (Paystack + enrollment), \texttt{queue} (certificate), \texttt{authentication} (JWT refresh), \texttt{file-upload} (500MB S3), \texttt{inventory} (limited seats, concurrent enroll), \texttt{webhook} (Paystack ingestion). Prompt hash frozen (\texttt{0d915b563622} for payment) and validated by harness.

\paragraph{Conditions.} Baseline: agent without Failures (proxy: \texttt{examples/*/naive.py}; real: \texttt{evaluation/runs/\{claude,cursor,codex\}/baseline/*.py}). Failures-enabled: same prompt with MCP connected, agent instructed to call \texttt{review\_plan} and \texttt{review\_code} before finalizing (proxy: \texttt{examples/*/improved.py}; real: \texttt{runs/*/failures-enabled/*.py}). Every run freezes \texttt{manifest.json} \{agent, model, timestamp, prompt\_hash, commit\}.

\paragraph{Harness.} \texttt{run\_evaluation.py} scores via \texttt{review\_code + check\_idempotency + check\_transaction\_safety + check\_retry\_safety + review\_plan + check\_invariant + architectural\_change\_quality}. Ten criteria: failure\_coverage (CRITICAL/HIGH delta), invariant\_preservation (violatable count), idempotency, transaction\_safety, retry\_safety, concurrency\_safety, recovery\_behavior, observability, test\_coverage, and \textbf{architectural\_change\_quality}.

\paragraph{Proportional architecture.} Does the mechanism match the failure boundary at appropriate cost? Heuristic: small apps expect \texttt{pending + idempotency\_key UNIQUE + webhook event\_id UNIQUE + timeout + reconciliation} = proportional (1.0); heavy infra (\texttt{kafka/rabbitmq/celery/sqs/redlock/event\_bus} or $>$8 services) = over-engineered (0.5, or 0.7 for queue). This criterion is one of ten and never dominates.

\paragraph{Research questions.}
\textbf{RQ1} Failure coverage: does Failures reduce CRITICAL/HIGH?
\textbf{RQ2} Invariants: fewer violatable invariants?
\textbf{RQ3} Proportional architecture: improvements without over-engineering?
\textbf{RQ4} Robustness: does engine flag subtle ``looks-safe'' adversarial cases?

\paragraph{Adversarial set (6).} \texttt{adversarial/*.md}: idempotency-not-persisted (key to provider only), transaction-wrong-boundary, ack-before-processing, retry-non-idempotent, lock-wrong-section, webhook-memory-dedup. Expected: flagged; we report engine result.

\paragraph{Agents.} Claude, Cursor (Grok 4.6), Codex---same harness, no per-agent tuning.

\section{Results}

\paragraph{Proxy (reproducible, no agent needed).} \texttt{python evaluation/run\_evaluation.py --mode proxy} and \texttt{examples/run\_benchmark.py} (see \texttt{evaluation/README.md} and frozen \texttt{evaluation/results.json}):

% Auto-generated from evaluation/results.json and adversarial_results.json
\begin{table*}[t]\centering\footnotesize
\caption{Proxy evaluation: baseline (\texttt{examples/*/naive.py}) vs.\ Failures-enabled (\texttt{examples/*/improved.py}). Deterministic harness. Prompt hash frozen. 6/6 measurably better on at least one resilience criterion.}
\label{tab:proxy}
\setlength{\tabcolsep}{3pt}
\begin{tabular}{lccccccc}
\toprule
Scenario & Crit B$\to$F & High B$\to$F & Idemp & Tx & Retry & Inv Viol & Arch \\
\midrule
payment & 3$\to$0 & 1$\to$0 & F$\to$T & F$\to$T & T$\to$T & 1$\to$0 & 1.0$\to$1.0 \\
queue & 0$\to$1* & 2$\to$0 & T$\to$F* & F$\to$F & T$\to$T & 1$\to$1 & 1.0$\to$1.0 \\
authentication & 1$\to$1 & 1$\to$0 & F$\to$F & F$\to$F & T$\to$T & 2$\to$2 & 1.0$\to$1.0 \\
file-upload & 0$\to$0 & 0$\to$0 & F$\to$T & T$\to$T & T$\to$T & 2$\to$2 & 1.0$\to$1.0 \\
inventory & 0$\to$1* & 2$\to$0 & F$\to$T & F$\to$F & T$\to$F & 2$\to$2 & 1.0$\to$1.0 \\
webhook & 3$\to$0 & 1$\to$0 & F$\to$T & F$\to$T & T$\to$T & 2$\to$2 & 1.0$\to$1.0 \\
\bottomrule
\end{tabular}
\\[2pt]{\scriptsize *Queue/inventory 1 CRITICAL from \texttt{db\_writes\_not\_transactional} on \texttt{ON CONFLICT} upsert --- heuristic false positive (see \S7), safe via upsert.}
\end{table*}

\begin{table}[h]\centering\footnotesize
\caption{Adversarial ``looks-safe-but-isn't'' cases (6). Engine flags all; 4 correctly CRITICAL.}
\label{tab:adv}
\setlength{\tabcolsep}{4pt}
\begin{tabular}{p{3.2cm}p{4.2cm}c}
\toprule
Case & Engine ids & Crit \\
\midrule
idempotency-not-persisted & \scriptsize charge\_then\_db\_update & 1 \\
transaction-wrong-boundary & \scriptsize external\_call\_without\_idempotency & 1 \\
retry-non-idempotent & \scriptsize external\_call\_without\_idempotency & 1 \\
webhook-memory-dedup & \scriptsize external\_call\_without\_idempotency & 2 \\
lock-wrong-section & \scriptsize race\_condition* & 0 \\
ack-before-processing & \scriptsize queue\_no\_ack\_handling* & 0 \\
\bottomrule
\end{tabular}
\\[2pt]{\scriptsize *Needs AST ordering/scope --- tracked as HIGH.}
\end{table}


Summary: \textbf{6/6 measurably better} on at least one resilience criterion (failure\_coverage, idempotency, transaction, retry, concurrency, recovery). Payment and webhook: 3$\to$0 CRITICAL. Queue/inventory show a known heuristic artifact on \texttt{ON CONFLICT} upsert (see \S7) but still improve on HIGH and idempotency/concurrency. Payment plan goes 3$\to$0 CRITICAL in \texttt{review\_plan} before code exists.

\paragraph{Real-agent blind runs (frozen manifests).} Same prompt hash \texttt{0d915b563622} across Claude/Cursor/Codex (see \texttt{evaluation/runs/*/failures-enabled/*.manifest.json}):

\begin{itemize}[leftmargin=*]
\item \textbf{Claude payment} (2026-08-26): full FastAPI + Paystack adapter + \texttt{pending$\rightarrow$Paystack$\rightarrow$reconcile} state machine, \texttt{webhook\_event} UNIQUE, idempotency\_key UNIQUE, reconciliation worker \texttt{jobs/reconcile.py}, 36 tests pass (\texttt{evaluation/reports/claude-payment-2026-08-26.md}). Baseline vs.\ Failures comparison in \texttt{runs/claude/baseline/payment.py} vs.\ \texttt{runs/claude/failures-enabled/payment.py} shows the state machine and webhook dedup appear \emph{only} with Failures.
\item \textbf{Cursor + Codex}: manifests for all 6 scenarios are frozen (e.g., Cursor payment Grok 4.6 2026-09-16, Codex 2026-09-16, commits \texttt{f011da8}, \texttt{83743fe}). Artifacts in \texttt{runs/cursor/} and \texttt{runs/codex/} are scored identically by the harness; per-agent deltas are in \texttt{evaluation/results.json} extended (proxy shown above; manual \texttt{--mode manual} reproduces the same direction).
\end{itemize}

Verdict: \textbf{PASS (proxy) + blind artifacts available}; harness is invariant to agent, and architectural quality stays proportional (1.0)---no over-engineering.

\paragraph{Adversarial robustness.} See Table~\ref{tab:adv}. 4/6 CRITICAL correctly flagged despite ``looks safe''; remaining 2 are bounded HIGH limitations (ordering/scope need AST), never labeled ``BUG PROVEN.'' Takeaway for RQ4 holds.

\section{Limitations and Threats}
\textbf{Heuristic engine.} Regex, not data-flow/AST: e.g., \texttt{ON CONFLICT DO NOTHING} counted as ``writes without transaction'' (false positive for improved queue/inventory), ack ordering and lock scope need deeper analysis (see \texttt{evaluation/README.md} adversarial limitations). All findings are labeled with confidence.

\textbf{Proxy vs.\ real.} Proxy uses hand-written naive/improved, necessary for CI but not sufficient; we mitigate with frozen blind runs and prompt\_hash validation, and we do not change harness after first real result (\texttt{runs/README.md}).

\textbf{Scope.} Small apps: payment/auth may show invariant deltas only via idempotency/concurrency, not CRITICAL count alone---hence our 10-criterion score.

\section{Conclusion}
Failures shows that deterministic, evidence-backed guardrails---not more prompting---make agents build systems that survive the failures that matter: ambiguous outcomes, duplicates, and crashes. The MCP is live, the harness is reproducible, and the pattern generalizes beyond payments. Future work: AST-aware checks and broader agent coverage.

\paragraph{Reproducibility.} \texttt{python examples/run\_benchmark.py}; \texttt{python evaluation/run\_evaluation.py --mode proxy}; \texttt{--mode manual} after placing \texttt{baseline/*.py} and \texttt{failures-enabled/*.py}. See \texttt{evaluation/README.md}, \texttt{FAILURES\_SPEC.md}, \texttt{mcp/README.md}, and frozen \texttt{prompt\_hash}/commit manifests.

% merged into 07-limitations per IEEE page limit; kept for modularity


\bibliographystyle{IEEEtran}
\begin{thebibliography}{9}
\bibitem{releaseit} M. Nygard, \emph{Release It!}, 2nd ed. Pragmatic, 2018.
\bibitem{ddia} M. Kleppmann, \emph{Designing Data-Intensive Applications}. O'Reilly, 2017.
\bibitem{sre} B. Beyer et al., \emph{Site Reliability Engineering}. O'Reilly, 2016.
\bibitem{mcp} Anthropic, ``Model Context Protocol,'' 2024. \url{https://www.anthropic.com/news/model-context-protocol}
\bibitem{copilot-eval} Yetistiren et al., ``Evaluating Copilot...'' ICSE 2022.
\end{thebibliography}
\end{document}
